\sect{Фундамент обучения с подкреплением}{rl-base}
%
В настоящем разделе кратко изложен математический аппарат
обучения с подкреплением в объёме, необходимом для разбора
работ, применяющих RL к задаче музыкального
аккомпанемента~\cite{Dorfer2018RL,Jiang2020RLDuet,Peter2023,
Wu2024ReaLchords}. Основная справочная литература~---
учебник Саттона и Барто~\cite{SuttonBarto2018}; для частично
наблюдаемых задач существенна работа
Кэлблинга-Литмана-Кассандры~\cite{Kaelbling1998POMDP}.

\paragraph*{Марковский процесс принятия решений.} Дискретный
марковский процесс принятия решений (MDP) задаётся пятёркой
$\mathcal{M}=(\mathcal{S},\mathcal{A},\mathcal{T},r,\gamma)$,
где $\mathcal{S}$~--- множество состояний, $\mathcal{A}$~---
множество действий, $\mathcal{T}(s'|s,a)$~--- ядро переходов
$\Prb(s_{t+1}=s'|s_t=s,a_t=a)$, $r(s,a)$~--- функция
вознаграждения и $\gamma\in[0,1]$~--- дисконт-фактор. Стратегия
агента $\pi(a|s)$ задаёт условное распределение действий по
состояниям; в случае стохастической стратегии~--- собственно
распределение, в случае детерминированной~---
$\pi:\mathcal{S}\to\mathcal{A}$. Цель агента~--- максимизация
ожидаемой суммарной дисконтированной награды
\begin{equation}
J(\pi) = \Expect_\pi\biggl[\sum_{t=0}^{\infty}\gamma^t r(s_t, a_t)\biggr].
\eqlab{rl-objective}
\end{equation}
Стандартные ценностные функции:
$V^\pi(s) = \Expect_\pi[\sum_{t=0}^\infty \gamma^t r_t \mid s_0=s]$~---
функция ценности состояния под политикой~$\pi$;
$Q^\pi(s,a) = \Expect_\pi[\sum_{t=0}^\infty \gamma^t r_t \mid s_0=s,
a_0=a]$~--- функция ценности пары состояние-действие;
\textit{функция преимущества} $A^\pi(s,a)=Q^\pi(s,a)-V^\pi(s)$
характеризует относительную ценность действия~$a$ в
состоянии~$s$~\cite[гл.~3]{SuttonBarto2018}. Ценностные функции
удовлетворяют уравнениям Беллмана; в частности,
$V^\pi(s) = \Expect_{a\sim\pi}\bigl[r(s,a) +
\gamma\Expect_{s'}V^\pi(s')\bigr]$.

\paragraph*{Частично наблюдаемые задачи (POMDP).} В реальных
условиях агент не наблюдает истинное состояние~$s$, а получает
лишь зашумлённое наблюдение $o\sim O(\cdot|s)$. Формализм
POMDP~\cite{Kaelbling1998POMDP} расширяет MDP добавлением
пространства наблюдений и распределения $O$. Стандартный приём
сведения POMDP к MDP~--- замена истинного состояния на
\textit{belief-состояние} $b_t = \Prb(s_t \mid o_{1:t}, a_{1:t-1})$,
которое обновляется байесовски при получении новых
наблюдений. В контексте отслеживания партитуры роль belief
играет нормированная forward-переменная $\tilde\alpha_t$ из
\Eqref{forward-rec}.

\paragraph*{Policy gradient и теорема о градиенте политики.}
Прямой оптимизации $J(\pi_\theta)$ по параметрам~$\theta$
посвящено семейство методов \textit{policy gradient}.
Теорема о градиенте политики гласит:
\begin{equation}
\nabla_\theta J(\pi_\theta) = \Expect_{\pi_\theta}\!\biggl[
\sum_{t=0}^{\infty} \nabla_\theta\log\pi_\theta(a_t|s_t)\,
A^{\pi_\theta}(s_t, a_t)
\biggr].
\eqlab{policy-gradient}
\end{equation}
В простейшей реализации REINFORCE
вместо~$A^\pi$ используется суммарная награда $G_t=\sum_{k\ge 0}
\gamma^k r_{t+k}$~\cite[гл.~13]{SuttonBarto2018}; в
варианте с baseline~--- разность $G_t - V_\phi(s_t)$, где
$V_\phi$~--- параметризованная оценка ценности. Дорфер с
коллегами~\cite{Dorfer2018RL} в задаче отслеживания партитуры
по картинке используют именно эту схему с дополнительной
оптимизацией advantage actor-critic (A2C) с
шестнадцатью параллельными акторами и обновлением каждые
$t_{\max}=15$ шагов.

\paragraph*{Generalized Advantage Estimation.} Шульман и
коллеги~\cite{Schulman2016GAE} предложили компромисс между
смещением и дисперсией оценки $A^\pi$ через экспоненциально
взвешенную сумму $k$-шаговых временных
разностей. С учётом обозначения
$\delta_t^V = r_t + \gamma V(s_{t+1}) - V(s_t)$
оценка GAE имеет вид~\cite[уравн.~16]{Schulman2016GAE}:
\begin{equation}
\hat A_t^{\,\text{GAE}(\gamma,\lambda)}
= \sum_{l=0}^{\infty} (\gamma\lambda)^l \,\delta_{t+l}^V.
\eqlab{gae-def}
\end{equation}
Параметры $\gamma$ и $\lambda$ управляют торгом между смещением
($\lambda=0$~--- минимальное смещение, максимальная дисперсия)
и дисперсией ($\lambda=1$~--- Монте-Карло). Стандартный диапазон
для практических задач~--- $\lambda\in[0.92, 0.98]$~\cite{Schulman2016GAE}.

\paragraph*{Proximal Policy Optimization.} Шульман и
коллеги~\cite{Schulman2017PPO} предложили
вычислительно эффективный заменитель Trust Region Policy
Optimization, в котором ограничение KL-дивергенции заменяется
явным «обрезанным» (clipped) суррогатным
функционалом. С учётом обозначения отношения
правдоподобий $r_t(\theta) = \pi_\theta(a_t|s_t)\bigl/
\pi_{\theta_{\text{old}}}(a_t|s_t)$ цель PPO имеет
вид~\cite[уравн.~7]{Schulman2017PPO}:
\begin{equation}
L^{\text{CLIP}}(\theta)
= \Expect_t\!\Bigl[\,\min\bigl(r_t(\theta)\,\hat A_t,\;
\mathrm{clip}(r_t(\theta), 1-\epsilon, 1+\epsilon)\,\hat A_t\bigr)
\,\Bigr],
\eqlab{ppo-clip}
\end{equation}
где $\hat A_t$~--- оценка преимущества (как правило, GAE),
а $\epsilon\in(0,1)$~--- параметр обрезки (типично $\epsilon=0.2$).
Функционал максимизируется методом мини-батч SGD с
несколькими эпохами по одному и тому же набору переходов.
Полная функция потерь обычно дополняется квадратичной потерей
по value-функции и поощрением энтропии:
$L_t(\theta,\phi) = L_t^{\text{CLIP}}(\theta) - c_1
\bigl(V_\phi(s_t) - V_t^{\text{target}}\bigr)^2
+ c_2 \mathcal{H}[\pi_\theta(\cdot|s_t)]$
с гиперпараметрами $c_1, c_2$.

\paragraph*{Behavioural cloning.} В случаях, когда доступна
демонстрация~--- последовательность пар (состояние, желаемое
действие)~--- начальная политика может быть обучена
супервизированно. Цель имеет вид
\begin{equation}
\mathcal{L}_{\text{BC}}(\theta) = \Expect_{(s,a^*)\sim\mathcal{D}}
\bigl\|\pi_\theta(s) - a^*\bigr\|^2 \quad \text{(непрерывный случай)}
\eqlab{bc-loss}
\end{equation}
или, в дискретном случае, кросс-энтропии. Behavioural cloning
обычно используется как этап инициализации перед последующей
RL-настройкой~\cite{Wu2024ReaLchords,Jaques2017SeqTutor}; в
чистом виде он подвержен \textit{distributional shift}, поскольку
при отклонениях от демонстрационного распределения политика
оказывается в незнакомых состояниях.

\paragraph*{Обучение по предпочтениям (RLHF).} Когда явная
функция вознаграждения недоступна, но доступны парные сравнения
коротких сегментов траектории, применяется схема обучения
по предпочтениям~\cite{Christiano2017Preference}. Вероятность
предпочтения сегмента $\sigma_1$ сегменту $\sigma_2$ моделируется
по схеме Брэдли-Терри~\cite{BradleyTerry1952}:
\begin{equation}
\hat\Prb[\sigma_1 \succ \sigma_2] = \frac{\exp\bigl[\sum_t
\hat r(o_t^1, a_t^1)\bigr]}{\exp\bigl[\sum_t \hat r(o_t^1, a_t^1)
\bigr] + \exp\bigl[\sum_t \hat r(o_t^2, a_t^2)\bigr]},
\eqlab{rlhf-pref}
\end{equation}
а функция вознаграждения~$\hat r$ обучается минимизацией
кросс-энтропии между этими предсказаниями и фактическими
пометками экспертов~\cite[уравн.~1]{Christiano2017Preference}.
В контексте музыкального применения RLHF может быть полезна
для обучения «эстетических» аспектов исполнения, не сводимых
к одной численной метрике (см.~\Secref{discussion}).

\paragraph*{KL-контроль.} Жак с коллегами~\cite{Jaques2017SeqTutor}
рассмотрели задачу дообучения предобученной модели
последовательности (например, MLE-обученной LSTM) с эвристическим
внешним вознаграждением, не разрушая статистических свойств
исходного распределения. Цель имеет вид
\begin{equation}
\max_\theta \Expect_{\tau\sim\pi_\theta}\!\bigl[r(\tau)\bigr]
- \beta\,\mathrm{KL}\bigl(\pi_\theta\,\|\,\pi_{\text{prior}}\bigr),
\eqlab{kl-control}
\end{equation}
где $\pi_{\text{prior}}$~--- замороженная исходная политика, а
$\beta>0$~--- регулирующий параметр. Тот же приём в современной
литературе известен как \textit{conservative fine-tuning} и
лежит в основе подхода ReaLchords для онлайн-генерации
аккомпанемента~\cite{Wu2024ReaLchords}.
